\section{解决方案架构：双 Agent 协作模式}

Anthropic 提出的解决方案包含两个关键角色：Initializer Agent 和 Coding Agent。两者通过共享的文件系统状态（特别是 progress 文件和 git history）实现跨 session 的信息传递。

\subsection{Initializer Agent：环境搭建}

Initializer Agent 仅在项目的第一个 session 中运行，负责创建后续所有 Coding Agent 所需的环境。它的具体职责包括：

\begin{enumerate}
    \item 创建 \texttt{init.sh} 脚本——用于启动开发服务器等基础设施
    \item 创建 \texttt{claude-progress.txt} 文件——记录各 Agent session 的工作日志
    \item 生成初始 git commit——记录所有被添加的文件
    \item 编写 feature list（功能清单）——基于用户 prompt 扩展出完整的功能需求列表
\end{enumerate}

\subsection{Coding Agent：增量式推进}

每个后续的 session 都运行 Coding Agent，其行为模式为：

\begin{enumerate}
    \item 阅读 progress 文件和 git log，了解当前进展
    \item 从 feature list 中选择下一个待实现的功能
    \item 实现该功能并进行测试
    \item 提交 git commit 并更新 progress 文件
    \item 确保环境处于干净状态后结束 session
\end{enumerate}

\begin{knowledgebox}{为什么用两个不同的 Agent？}
使用不同的 prompt 区分"首次运行"和"后续运行"是一个重要的设计选择。Initializer Agent 的 prompt 侧重于环境搭建和需求分析，而 Coding Agent 的 prompt 侧重于增量开发和代码质量。这种分工避免了单一 prompt 过于复杂，也让每个 Agent 能更专注于自己的职责。Anthropic 的 Claude 4 prompting guide 中也推荐了这种"第一个 context window 使用不同 prompt"的 harness 结构。
\end{knowledgebox}

\subsection{信息传递机制}

两个 Agent 之间的信息传递完全依赖文件系统，而非任何内存或数据库：

\begin{table}[H]
\centering
\begin{tabular}{p{4cm}p{4cm}p{5cm}}
\toprule
\textbf{信息载体} & \textbf{作用} & \textbf{设计要点} \\
\midrule
\texttt{claude-progress.txt} & 记录每个 session 的工作摘要 & 自由文本格式，方便 Agent 快速理解上下文 \\
\midrule
\texttt{feature\_list.json} & 追踪所有功能的完成状态 & JSON 格式，防止 Agent 随意修改内容 \\
\midrule
Git history & 记录所有代码变更 & 描述性的 commit message，支持 revert 操作 \\
\midrule
\texttt{init.sh} & 环境启动脚本 & Agent 无需重新发现如何运行项目 \\
\bottomrule
\end{tabular}
\caption{跨 session 信息传递的四种载体}
\end{table}

\subsection{本章小结}

双 Agent 架构将"一次性完成"的不切实际目标分解为"搭建环境 + 持续增量推进"。Initializer Agent 负责创建所有必要的基础设施，Coding Agent 在每个 session 中做出可管理的增量进展。两者之间通过 progress 文件、feature list、git history 和启动脚本实现状态传递。


